\chapter*{\hfill{\centering ВВЕДЕНИЕ}\hfill}

\addcontentsline{toc}{chapter}{ВВЕДЕНИЕ}

Тема поиска всегда остаётся в центре внимания в информационных технологиях, задача поиска информации в структурах данных является одной из базовых задач программирования. Оптимизация поиска~---~одна из важных задач, её решение позволяет ускорить работу различных баз и структур данных~\cite{book_knut}.

Целью данной работы является исследование методов бинарного поиска и поиска полным перебором в массиве, а также анализ этих алгоритмов по количеству сравнений элементов.

\label{sec:goals}
Для достижения поставленной цели необходимо выполнить следующие задачи:
\begin{itemize}[label=---]
	\item описать линейный и бинарный алгоритмы поиска;
	\item написать программу, реализующую алгоритм поиска полным перебором и бинарный алгоритм поиска;
	\item проверить работоспособность программы, выполнив модульное тестирование;
	\item замерить число сравнений элементов в упомянутых выше алгоритмах при разном размере массива и положении искомого элемента;
	\item провести анализ полученных результатов.
\end{itemize}




